% Import the fullpage package to adjust page margins
\usepackage{fullpage}
\usepackage{tikz}
\usepackage{pgfplots}
\pgfplotsset{compat=1.18}
\usepackage{tikz,pgfplots,caption}
\usepackage{graphicx}
\usepackage{adjustbox}
\usepackage{multicol}
\usepackage{tabularx}
\usepackage{array}
\usepackage{needspace}
\usepackage{tikz}
\usetikzlibrary{shapes,decorations,arrows,calc,arrows.meta,fit,positioning}
\newcolumntype{L}{>{\raggedright\arraybackslash}X}
\usepackage{longtable}
\usepackage{dcolumn}
\usepackage{xurl} % allows urls to break nicely across lines (recommended)

% Import the enumitem package to customize lists
\usepackage{enumitem}

% Import various AMS (American Mathematical Society) packages for mathematical symbols and formatting
\usepackage{amsfonts, amssymb, amsmath, amsthm}

% Import the mathtools package for advanced math typesetting
\usepackage{mathtools}

%For tables
\usepackage{multirow}

% Configure hyperlinks in the PDF file (e.g., author name, title)
\usepackage[pdftex, pdfauthor={\name}, pdftitle={\classnum~\assignment}]{hyperref}

% Import the xcolor package with the 'dvipsnames' option to define additional colors
\usepackage[dvipsnames]{xcolor}

% Import additional math fonts
\usepackage{bbm}

% Allow inserting images into the document
\usepackage{graphicx}

% Import additional mathematical script fonts
\usepackage{mathrsfs}

% Allow including entire PDF pages
\usepackage{pdfpages}

% Enable flexible table column formatting
\usepackage{tabularx}

% Enable landscape orientation for pages
\usepackage{pdflscape}

% Improve table formatting with multi-line cells
\usepackage{makecell}

% Improve table design and formatting
\usepackage{booktabs}

% Import natbib for better citation handling
\usepackage{natbib}

% Customize captions for tables and figures
\usepackage{caption}

% Allow subfigures within figures
\usepackage{subcaption}

% Import physics package for easier notation of derivatives, vectors, etc.
\usepackage{physics}

% Import the tcolorbox package to create colored boxes (e.g., for solutions or notes)
\usepackage[many]{tcolorbox}

% Enable conditional text inclusion (e.g., for draft versions)
\usepackage{version}

% Enable conditional logic in LaTeX (used for macros)
\usepackage{ifthen}

% Strike through mathematical expressions
\usepackage{cancel}

% Import the listings package to format source code
\usepackage{listings}

% Use the Courier font for code listings
\usepackage{courier}

% Import TikZ for creating diagrams and illustrations
\usepackage{tikz}

% Import a package for drawing game theory diagrams
\usepackage{istgame}

% Allow better placement of floating figures and tables
\usepackage{float}

% Customize hyperlink colors (e.g., blue links)
\hypersetup{
	colorlinks=true,
	linkcolor=blue,
	filecolor=magenta,
	urlcolor=blue,
}

% Adjust paragraph formatting
\setlength{\parindent}{0mm}  % No indentation at the start of paragraphs
\setlength{\parskip}{2mm}    % Space between paragraphs

% Customize numbered lists: (a), (b), (c), etc.
\setlist[enumerate]{label=({\alph*})}

% Customize nested numbered lists: (i), (ii), (iii), etc.
\setlist[enumerate, 2]{label=({\roman*})}

% Allow equations to break across pages (level 1 means minimal breaks)
\allowdisplaybreaks[1]

% Define a macro to generate a problem set header
\newcommand{\psetheader}{
	\ifthenelse{\isundefined{\collaborators}}{ % Check if collaborators are specified
		\begin{center}
			{\setlength{\parindent}{0cm} \setlength{\parskip}{0mm}
				
				% Display course number, semester, and assignment title on the left, name on the right
				{\textbf{\classnum~\semester:~\assignment} \hfill \name}
				
				% Display subject and email, with email as a clickable link
				\subject \hfill \href{mailto:\email}{\tt \email}
				
				% Display instructor names and due date
				Instructor:~\instructors \hfill Due Date:~\duedate	
				
				% Horizontal rule (separator line)
				\hrulefill}
		\end{center}
	}{
		\begin{center}
			{\setlength{\parindent}{0cm} \setlength{\parskip}{0mm}
				
				% If collaborators exist, add them as a footnote
				{\textbf{\classnum~\semester:~\assignment} \hfill \name\footnote{Collaborator(s): \collaborators}}
				
				% Display subject and email
				\subject \hfill \href{mailto:\email}{\tt \email}
				
				% Display instructor names and due date
				Instructor(s):~\instructors \hfill Due Date:~\duedate	
				
				\hrulefill}
		\end{center}
	}
}

% Customize page numbering: include class number, assignment, and page number
\usepackage{fancyhdr}
\pagestyle{fancy}
\fancyhf{}
\lfoot{\classnum~\assignment}
\rfoot{\thepage}
\renewcommand{\headrulewidth}{0pt}
\renewcommand{\footrulewidth}{0pt}

% Define a macro to format point values for problems
\makeatletter
\def\points{\@ifnextchar[{\@with}{\@without}}
\def\@with[#1]#2{{\ifthenelse{\equal{#2}{1}}{{[1 point, #1]}}{{[#2 points, #1]}}}}
\def\@without#1{\ifthenelse{\equal{#1}{1}}{{[1 point]}}{{[#1 points]}}}
\makeatother

% ---------------------------------------
% Elegant Theorem Style (Academic / Journal)
% ---------------------------------------
\newcommand{\MRS}{\operatorname{MRS}}

\makeatletter
\newtheoremstyle{econ-thm}% name
  {6pt}% Space above
  {6pt}% Space below
  {\normalfont}% Body font
  {}% Indent
  {\bfseries}% Default head font (BOLD ONLY)
  {.}% Punctuation after theorem number
  {0.5em}% Space after theorem head
  {%
    \textbf{\thmname{#1}\thmnumber{ #2}.}% <-- LABEL: bold italic
    \ifx\thmnote\@empty
    \else\ \textbf{\thmnote{#3}}\fi% <-- NAME: bold only
  }% Head spec
\makeatother

\theoremstyle{econ-thm}
\newtheorem{theorem}{Theorem}
\newtheorem{lemma}{Lemma}
\newtheorem{corollary}{Corollary}
\newtheorem{proposition}{Proposition}
\newtheorem{definition}{Definition}



% Define an environment for problems (sets text color to black)
\newenvironment{problem}[1]{\color{black} #1}{}

% Define a solution box with a green border and automatic proof environment
\newenvironment{solution}{%
  \leavevmode
  \begin{tcolorbox}[
    breakable,
    colback=green!5!white,
    colframe=green!75!black,
    enhanced jigsaw,
    % ↓ ensure no first-line indent anywhere in the box
    before upper={\parindent0pt\relax \parskip2mm}%
  ]%
  \proof[\scshape Solution:] \mbox{}\par\noindent
}{%
  \renewcommand{\qedsymbol}{$\blacksquare$}%
  \endproof
  \end{tcolorbox}%
}

% Define a reflection environment with a gray box
\newenvironment{reflection}{\begin{tcolorbox}[breakable, colback=black!8!white,colframe=black!60!white, enhanced jigsaw, parbox = false]\textsc{Reflections:}}{\end{tcolorbox}}

% Define a custom QED symbol for proofs
\newcommand{\qedh}{\renewcommand{\qedsymbol}{$\blacksquare$}\qedhere}

% Define custom colors for syntax highlighting
\definecolor{mygreen}{rgb}{0,0.6,0}
\definecolor{mygray}{rgb}{0.5,0.5,0.5}
\definecolor{mymauve}{rgb}{0.58,0,0.82}

% Define Stata programming language syntax highlighting for listings
\lstdefinelanguage{Stata}{
	% Define categories of keywords for Stata commands
	morekeywords=[1]{regress, reg, summarize, sum, display, di, generate, gen, bysort, use, import, delimited, predict, quietly, probit, margins, test},
	% Define reserved words
	morekeywords=[2]{boolean, break, byte, case, catch, class, const, continue, delete, double, else, enum, export, float, function, if, int, local, long, matrix, namespace, new, numeric, NULL, operator, pointer, private, protected, public, return, scalar, short, static, string, struct, super, switch, template, this, throw, typedef, typename, union, unsigned, using, vector, virtual, void, volatile, while,},
	% Define comment styles for Stata code
	morecomment=[l]{//},
	morecomment=[s]{/*}{*/},
	sensitive=true,
}

% Set global options for code listings (e.g., font, line numbering)
\lstset{ 
	backgroundcolor=\color{white},  
	basicstyle=\footnotesize\ttfamily, 
	numbers=left,
	numberstyle=\tiny\color{mygray},  
	keywordstyle=\color{blue},  
	commentstyle=\color{mygreen},  
	stringstyle=\color{mymauve},  
	frame=single,  
	tabsize=2,
}
